\chapter{Технологический раздел}

Технологический раздел посвящен практической реализации гибридного метода восстановления глубины из монокулярных изображений.
В данном разделе рассматриваются ключевые аспекты программной реализации и тестирования разработанного решения.

Особое внимание уделяется обоснованию выбора программных средств и библиотек, необходимых для эффективной реализации алгоритмов.

Для объективной оценки качества реконструкции будут формализованы метрики: MAE, RMSE, REL и точность δ, позволяющие количественно оценивать эффективность метода на разных типах изображений.

\section{Выбор программных средств}

Для реализации методов восстановления глубины из монокулярных изображений был выбран язык программирования Python.
Этот выбор обусловлен наличием в его экосистеме всех необходимых библиотек и инструментов для комплексного решения поставленной задачи.

% TODO добавить ссылки на источники

\begin{enumerate}
    \item \textbf{Библиотеки обработки изображений:} OpenCV~\cite{opencv_library}, scikit-image~\cite{van2014scikit} и PIL/Pillow~\cite{clark2015pillow} — обеспечивают функционал для эффективной предобработки архивных фотографий, включая фильтрацию, сегментацию и работу с различными форматами файлов.
    \item \textbf{Библиотеки машинного обучения:} PyTorch~\cite{paszke2019pytorch}, TensorFlow/Keras~\cite{abadi2016tensorflow} и scikit-learn~\cite{pedregosa2011scikit} — предоставляют инструменты для применения предобученной модели MiDaS для оценки глубины, что исключает необходимость обучения собственной нейронной сети и повышает эффективность разработки.
    \item \textbf{Библиотеки научных вычислений:} NumPy~\cite{harris2020array}, SciPy~\cite{virtanen2020scipy}, Matplotlib~\cite{hunter2007matplotlib} и Open3D~\cite{zhou2018open3d} — обеспечивают эффективную работу с многомерными массивами, методы оптимизации, визуализацию результатов и обработку трехмерных данных.
\end{enumerate}

Таким образом, совокупность доступных в Python библиотек делает его достаточным инструментом для решения всех аспектов задачи восстановления глубины, включая предобработку данных, использование готовой модели MiDaS и визуализацию результатов.

\section{Описание входных и выходных данных}

Входными данными для алгоритмов восстановления глубины являются монокулярные цветные фотографии в различных форматах (JPEG, PNG, TIFF).

Результатом работы алгоритмов восстановления глубины являются карта глубин в виде одноканальных изображений, где:

\begin{itemize}
    \item значение пикселя соответствует относительному расстоянию от камеры (меньшие значения — ближе к камере);
    \item диапазон значений нормализован от 0 до 1 (или от 0 до 255 для 8-битных изображений);
    \item разрешение карты глубины соответствует разрешению входного изображения.
\end{itemize}


\section{Реализация методов}

% TODO добавить Листингиииииииииииии

В листинге~\ref{lst:preproc} представлен код предварительной обработки входного изображения, включающий методы шумоподавления, нормализации освещения и коррекции артефактов, характерных для архивных фотографий.

\begin{center}
    \captionsetup{justification=raggedright,singlelinecheck=off}
    \begin{lstlisting}[label=lst:preproc,caption=Предварительная обработка изображения]
        import cv2
        import numpy as np


        def preprocess_archival_photo(image_path, brightness_boost=1.5, contrast_boost=1.2):
            if isinstance(image_path, str):
                image = cv2.imread(image_path)
                if image is not None and len(image.shape) == 3:
                    image = cv2.cvtColor(image, cv2.COLOR_BGR2RGB)
            else:
                image = image_path.copy()

            # Проверяем, является ли изображение цветным или оттенков серого
            if len(image.shape) == 3:
                gray = cv2.cvtColor(image, cv2.COLOR_RGB2GRAY)
            else:
                gray = image.copy()

            # Выявление и маскирование дефектов
            blur = cv2.GaussianBlur(gray, (5, 5), 0)
            diff = cv2.absdiff(gray, blur)
            _, artifact_mask = cv2.threshold(diff, 25, 255, cv2.THRESH_BINARY)

            kernel = np.ones((3, 3), np.uint8)
            artifact_mask = cv2.morphologyEx(artifact_mask, cv2.MORPH_CLOSE, kernel)
            artifact_mask = cv2.dilate(artifact_mask, kernel, iterations=1)

            # Инпейнтинг для удаления дефектов
            cleaned = cv2.inpaint(gray, artifact_mask.astype(np.uint8), 3, cv2.INPAINT_TELEA)

            # Повышение контраста
            clahe = cv2.createCLAHE(clipLimit=2.0, tileGridSize=(8, 8))
            enhanced = clahe.apply(cleaned)

            # Повышение яркости
            mean_brightness = np.mean(enhanced)
            target_brightness = 80

            brightness_factor = min((target_brightness / (mean_brightness + 1e-6)), brightness_boost)
            brightened = np.clip(enhanced * brightness_factor, 0, 255).astype(np.uint8)

            # Улучшение контраста
            alpha = min(contrast_boost, 1.1)
            beta = 128 * (1 - alpha)
            enhanced_contrast = np.clip(alpha * brightened + beta, 0, 255).astype(np.uint8)

            # Билатеральная фильтрация для сохранения краев при устранении шума
            denoised = cv2.bilateralFilter(enhanced_contrast, 7, 25, 25)

            # Нормализация до [0, 1]
            preprocessed = denoised.astype(np.float32) / 255.0

            # Гамма-коррекция для улучшения контраста в средних тонах
            gamma = 1.1
            preprocessed = np.power(preprocessed, gamma)

            return preprocessed, artifact_mask
    \end{lstlisting}
\end{center}

Листинг~\ref{lst:depth_extract} иллюстрирует основной алгоритм извлечения глубины, который объединяет нейросетевой метод для первичной оценки глубины с последующим уточнением деталей.

\begin{center}
    \captionsetup{justification=raggedright,singlelinecheck=off}
    \begin{lstlisting}[label=lst:depth_extract,caption=Извлечение глубины из подготовленного изображения]
        import os
import sys
import numpy as np
import cv2
import time

from scipy.ndimage import gaussian_filter, sobel

# Добавляем путь к директории code
sys.path.append(os.path.dirname(os.path.dirname(os.path.abspath(__file__))))

# Изменим импорты, чтобы они указывали на директории внутри code
from preprocessing.image_cleaner import preprocess_archival_photo
from preprocessing.segmentation import segment_by_color
from reconstruction.lighting import estimate_lighting
from reconstruction.albedo_shading import separate_albedo_shading
from reconstruction.normals import estimate_normals
from reconstruction.depth import recover_depth
from reconstruction.midas_depth import estimate_depth_midas
from postprocessing.refinement import postprocess_depth

class ArchivalDepthReconstruction:
    def __init__(self, iterations=200, lambda_smooth=0.15, regularization_weight=0.3,
                 brightness_boost=1.5, contrast_boost=1.2, n_color_segments=6,
                 method="segmentation", midas_model="MiDaS_small"):
        self.iterations = iterations
        self.lambda_smooth = lambda_smooth
        self.regularization_weight = regularization_weight
        self.brightness_boost = brightness_boost
        self.contrast_boost = contrast_boost
        self.n_color_segments = n_color_segments
        self.method = method
        self.midas_model = midas_model

    def process_hybrid(self, image_path):

        # Шаг 1: Получение базовой карты глубины
        print("Получение базовой карты глубины через MiDaS...")
        midas_results = self.process_midas(image_path)
        base_depth = midas_results['depth_map']
        preprocessed = midas_results['preprocessed']
        artifact_mask = np.zeros_like(preprocessed, dtype=np.uint8)

        # Шаг 2: Сегментация изображения
        print("Сегментация изображения для детализации...")
        segments_mask, bboxes = segment_by_color(
            preprocessed,
            self.n_color_segments
        )

        # Шаг 3: Инициализация карт деталей
        detail_depth = np.zeros_like(base_depth)
        detail_weight = np.zeros_like(base_depth)

        # Шаг 4: Цикл по сегментам
        print("Детализация каждого сегмента...")
        for segment_id, bbox in bboxes.items():
            y1, y2, x1, x2 = bbox

            segment_mask = (segments_mask[y1:y2, x1:x2] == segment_id).astype(np.float32)

            # Пропускаем слишком маленькие сегменты
            if np.sum(segment_mask) < 100:
                continue

            segment_img = preprocessed[y1:y2, x1:x2].copy()
            segment_base_depth = base_depth[y1:y2, x1:x2].copy()

            # Оценка условий освещения
            lighting_params = estimate_lighting(segment_img)

            # Разделение альбедо и затенения
            albedo, shading = separate_albedo_shading(segment_img, lighting_params)

            # Оценка нормалей поверхности
            normals, confidence_mask = estimate_normals(shading, lighting_params)

            # Восстановление деталей
            segment_details = recover_depth(
                normals,
                confidence_mask,
                iterations=self.iterations // 2,
                initial_depth=segment_base_depth
            )

            segment_weight = confidence_mask * np.abs(sobel(segment_img, axis=0)) * np.abs(sobel(segment_img, axis=1))
            segment_weight = gaussian_filter(segment_weight, sigma=1.0)
            if np.max(segment_weight) > 0:
                segment_weight = segment_weight / np.max(segment_weight)
            detail_depth[y1:y2, x1:x2] += segment_details * segment_mask * segment_weight
            detail_weight[y1:y2, x1:x2] += segment_mask * segment_weight

        # Шаг 5: Комбинация базовой карты и деталей
        detail_weight = np.maximum(detail_weight, 1e-6)
        detail_depth = detail_depth / detail_weight

        alpha_base = 0.7
        alpha_map = alpha_base - (detail_weight * 0.4)
        alpha_map = np.clip(alpha_map, 0.3, 0.9)

        final_depth = alpha_map * base_depth + (1 - alpha_map) * detail_depth


        # Формирование результатов
        results = {
            'preprocessed': preprocessed,
            'segments_mask': segments_mask,
            'depth_map': base_depth,
            'detail_depth': detail_depth,
            'detail_weight': detail_weight,
            'final_depth': final_depth,
            'bboxes': bboxes,
        }

        return results
    \end{lstlisting}
\end{center}

В листинге~\ref{lst:segmentation} представлен код сегментации изображения, позволяющий выделить различные объекты сцены для их дальнейшей раздельной обработки методом Shape From Shading.

\begin{center}
    \captionsetup{justification=raggedright,singlelinecheck=off}
    \begin{lstlisting}[label=lst:segmentation,caption=Сегментация изображения]
        import cv2
import numpy as np
from skimage.measure import regionprops
from sklearn.cluster import KMeans


def segment_by_color(img, n_color_segments=6):
    # Преобразование изображения в цветное
    if len(img.shape) == 2:
        image_for_cluster = np.stack([img] * 3, axis=2)
    else:
        image_for_cluster = img.copy()

    pixels = image_for_cluster.reshape(-1, 3)

    # Группировка пикселей по цветам
    kmeans = KMeans(n_clusters=n_color_segments, random_state=42, n_init=10)
    labels = kmeans.fit_predict(pixels)

    segmented = labels.reshape(img.shape[:2])

    h, w = segmented.shape[:2]
    bboxes = {}

    # Цикл по цветовым сегментам
    for label in range(n_color_segments):
        # Создание бинарной маски для текущего цвета
        mask = (segmented == label).astype(np.uint8)

        # Поиск связных компонентов
        num_labels, labels_img = cv2.connectedComponents(mask)

        for region_id in range(1, num_labels):
            region_mask = (labels_img == region_id).astype(np.uint8)

            # Получение прямоугольной ограничивающей рамки
            props = regionprops(region_mask)
            if not props:
                continue

            y1, x1, y2, x2 = props[0].bbox

            margin = 5
            y1 = max(0, y1 - margin)
            y2 = min(h, y2 + margin)
            x1 = max(0, x1 - margin)
            x2 = min(w, x2 + margin)

            # Если область слишком маленькая, она пропускается
            if (y2 - y1) < 20 or (x2 - x1) < 20:
                continue

            # Создание уникального идентификатора для этого прямоугольного сегмента
            segment_id = label * 1000 + region_id

            segmented[y1:y2, x1:x2][region_mask[y1:y2, x1:x2] > 0] = segment_id

            # Сохранение ограничивающей рамки
            bboxes[segment_id] = (y1, y2, x1, x2)

    return segmented, bboxes

    \end{lstlisting}
\end{center}

Листинг~\ref{lst:light_estimation} представляет реализацию алгоритма оценки параметров освещения сцены, необходимых для корректного применения метода Shape from Shading.

\begin{center}
    \captionsetup{justification=raggedright,singlelinecheck=off}
    \begin{lstlisting}[label=lst:light_estimation,caption=Оценка освещения]
import cv2
import numpy as np


def estimate_lighting(image):
    hist = cv2.calcHist([image * 255], [0], None, [256], [0, 256])
    hist_normalized = hist / (image.shape[0] * image.shape[1])
    brightness_mean = np.mean(image)
    brightness_std = np.std(image)

    # Вычисляем градиенты
    grad_x = cv2.Sobel(image, cv2.CV_32F, 1, 0, ksize=5)
    grad_y = cv2.Sobel(image, cv2.CV_32F, 0, 1, ksize=5)
    gradient_magnitude = np.sqrt(grad_x ** 2 + grad_y ** 2)

    # Определяем области с постоянным освещением
    smooth_regions = gradient_magnitude < np.percentile(gradient_magnitude, 25)

    # Оцениваем направление света по градиентам
    avg_grad_x = np.mean(grad_x)
    avg_grad_y = np.mean(grad_y)

    light_x = -avg_grad_x
    light_y = -avg_grad_y
    light_z = 1.0

    light_direction = np.array([light_x, light_y, light_z])
    norm = np.linalg.norm(light_direction)
    if norm > 0:
        light_direction = light_direction / norm

    lighting_params = {
        'direction': light_direction,
        'brightness_mean': brightness_mean,
        'brightness_std': brightness_std,
        'hist': hist_normalized,
        'smooth_regions': smooth_regions
    }

    return lighting_params

    \end{lstlisting}
\end{center}

В листинге~\ref{lst:albedo_separation} представлен код разделения альбедо и затенения метода Shape from Shading, позволяющий отделить собственные цвета поверхностей от эффектов освещения.

\begin{center}
    \captionsetup{justification=raggedright,singlelinecheck=off}
    \begin{lstlisting}[label=lst:albedo_separation,caption=Разделение альбедо и затенения]
# Содержимое файла albedo_separation.py
# Код разделения альбедо и затенения
    \end{lstlisting}
\end{center}

Листинг~\ref{lst:normals_estimation} содержит код вычисления нормалей к поверхности на основе полученного разделения альбедо и затенения.

\begin{center}
    \captionsetup{justification=raggedright,singlelinecheck=off}
    \begin{lstlisting}[label=lst:normals_estimation,caption=Оценка нормалей]
# Содержимое файла normals_estimation.py
# Код оценки нормалей
    \end{lstlisting}
\end{center}

Заключительный этап обработки представлен в листинге~\ref{lst:depth_integration} --- алгоритм интеграции поля нормалей для получения итоговой карты глубины с высоким уровнем детализации поверхностей.

\begin{center}
    \captionsetup{justification=raggedright,singlelinecheck=off}
    \begin{lstlisting}[label=lst:depth_integration,caption=Восстановление карты глубин по нормалям]
# Содержимое файла depth_integration.py
# Код восстановления карты глубин по нормалям
    \end{lstlisting}
\end{center}





\section{Метрики оценки качества реконструкции глубины}

% TODO добавить ссылки на источники

\subsection{Абсолютная средняя ошибка (MAE)}
\begin{equation}
\text{MAE} = \frac{1}{n} \sum_{i=1}^{n} |d_{i}^{\text{pred}} - d_{i}^{\text{gt}}|
\end{equation}

MAE представляет собой среднее абсолютное отклонение предсказанной глубины от эталонной. Эта метрика менее чувствительна к выбросам по сравнению с RMSE и позволяет выявлять систематические ошибки в предсказании глубины~\cite{zhang2025review}.

\subsection{Среднеквадратичная ошибка (RMSE)}
\begin{equation}
\text{RMSE} = \sqrt{\frac{1}{n} \sum_{i=1}^{n} (d_{i}^{\text{pred}} - d_{i}^{\text{gt}})^2}
\end{equation}

RMSE придаёт больший вес крупным ошибкам и является стандартной метрикой для оценки точности методов оценки глубины~\cite{vis3}. Благодаря своей чувствительности к выбросам, данная метрика эффективна при оценке методов, где критично минимизировать значительные отклонения.

\subsection{Относительная абсолютная ошибка (REL)}
\begin{equation}
\text{REL} = \frac{1}{n} \sum_{i=1}^{n} \frac{|d_{i}^{\text{pred}} - d_{i}^{\text{gt}}|}{d_{i}^{\text{gt}}}
\end{equation}

REL учитывает относительную величину ошибки, что особенно важно при работе с объектами, находящимися на разном расстоянии от камеры. Эта метрика предпочтительна при оценке качества реконструкции исторических объектов, где точная ш��ала глубины может быть неизвестна~\cite{zhang2025review}.

\subsection{Точность $\delta$}
\begin{equation}
\delta_{k} = \frac{1}{n} \sum_{i=1}^{n} \mathbb{I}\left[\max\left(\frac{d_{i}^{\text{pred}}}{d_{i}^{\text{gt}}}, \frac{d_{i}^{\text{gt}}}{d_{i}^{\text{pred}}}\right) < 1.25^k\right]
\end{equation}

где $k \in \{1, 2, 3\}$, а $\mathbb{I}[\cdot]$ — индикаторная функция, равная 1, если условие выполняется, и 0 в противном случае.

Эта метрика определяет процент пикселей, для которых отношение предсказанной глубины к истинной (или наоборот) не превышает заданного порога. Стандартно используются три порога: $1.25$, $1.25^2$ и $1.25^3$~\cite{vis2}.


\section{Тестирование}

Целью тестирования является верификация корректности функционирования разработанного гибридного метода восстановления глубины, а также оценка его устойчивости к различным эффектам деградации, характерным для исторических фотоматериалов.

Фундаментальной проблемой при проведении подобного тестирования выступает отсутствие достоверных эталонных данных о глубине для архивных фотографий, что исключает возможность прямой валидации точности реконструкции.

Для преодоления данного ограничения был сформирован комплексный методологический подход, основанный на использовании современных датасетов, снабженных верифицированными картами глубины. К изображениям из этих наборов применяются алгоритмы искусственного старения, моделирующие типичные дефекты исторических фотоматериалов. Данная методика обеспечивает возможность количественного сопоставления результатов реконструкции с эталонными данными и объективной оценки качества восстановления.

Предложенный подход позволяет проводить валидацию алгоритма в строго контролируемых условиях и идентифицировать потенциальные проблемные аспекты при обработке исторических фотографических материалов.

\subsection{Описание тестового набора}

Для проведения тестирования были выбраны три репрезентативных датасета, обеспечивающих разнообразие тестовых сцен.

% TODO добавить ссылки на источники датасетов и допустимых значений метрик

\begin{enumerate}
    \item \textbf{NYU Depth v2} --- 1449 изображений внутренних помещений с глубиной, полученной с помощью сенсора Kinect. Разрешение 640×480 пикселей~\cite{silberman2012indoor}.

    \item \textbf{Make3D} --- 534 изображения уличных сцен с глубиной, измеренной лазерным сканером. Разрешение 2272×1704 пикселей~\cite{saxena2009make3d}.

    \item \textbf{SceneNet RGB-D} --- 5M синтетических изображений с точными картами глубины. Разрешение 320×240 пикселей~\cite{mccormac2017scenenet}.
\end{enumerate}

Из каждого датасета было случайным образом отобрано по 50 изображений для тестирования, обеспечивая статистическую значимость результатов при сохранении выполнимости тестирования в разумные сроки.

Для имитации состаривания изображений был разработан программный модуль, реализующий следующие эффекты с параметрически задаваемой интенсивностью:

\begin{enumerate}
    \item \textbf{Сепия} — преобразование RGB в тонированное монохромное изображение (интенсивность 0.0-1.0).
    \item \textbf{Виньетирование} — затемнение краев изображения (радиус 0.5-0.9, интенсивность 0.2-0.8).
    \item \textbf{Шум и зернистость} — добавление шума, имитирующего фотопленку (плотность 0.01-0.5).
    \item \textbf{Выцветание} — снижение контраста и насыщенности (интенсивность 0.1-0.8).
    \item \textbf{Случайные пятна и царапины} — добавление локальных дефектов (плотность 0.01-0.3).
\end{enumerate}

\subsubsection{Классы эквивалентности}

\begin{enumerate}
    \item \textbf{Оригинальные изображения (без эффектов старения)}
    \begin{itemize}
        \item \textit{Тестовый набор:} по 50 изображений из каждого датасета
        \item \textit{Ожидаемые метрики:}
        \begin{itemize}
            \item MAE: 0.28-0.35
            \item RMSE: 0.55-0.70
            \item REL: 0.12-0.16
            \item $\delta_1$: 0.75-0.90
        \end{itemize}
    \end{itemize}

    \item \textbf{Изображения со средними эффектами старения (интенсивность 0.3-0.6)}
    \begin{itemize}
        \item \textit{Тестовый набор:} по 50 изображений из каждого датасета с умеренным состариванием
        \item \textit{Ожидаемые метрики:}
        \begin{itemize}
            \item MAE: 0.30-0.42
            \item RMSE: 0.61-0.88
            \item REL: 0.13-0.19
            \item $\delta_1$: 0.65-0.80
        \end{itemize}
    \end{itemize}

    \item \textbf{Изображения с сильными эффектами старения (интенсивность 0.6-0.9)}
    \begin{itemize}
        \item \textit{Тестовый набор:} по 50 изображений из каждого датасета с сильным состариванием
        \item \textit{Ожидаемые метрики:}
        \begin{itemize}
            \item MAE: 0.34-0.49
            \item RMSE: 0.69-1.02
            \item REL: 0.14-0.22
            \item $\delta_1$: 0.55-0.70
        \end{itemize}
    \end{itemize}

    \item \textbf{Негативный тест: экстремально низкое разрешение}
    \begin{itemize}
        \item \textit{Тестовый набор:} 50 изображений с разрешением менее 100×100 пикселей
        \item \textit{Ожидаемые метрики:}
        \begin{itemize}
            \item MAE: >0.51
            \item RMSE: >1.09
            \item REL: >0.22
            \item $\delta_1$: <0.50
        \end{itemize}
    \end{itemize}

    \item \textbf{Негативный тест: экстремальный уровень шума}
    \begin{itemize}
        \item \textit{Тестовый набор:} 50 изображений с SNR < 5 дБ
        \item \textit{Ожидаемые метрики:}
        \begin{itemize}
            \item MAE: >0.51
            \item RMSE: >1.09
            \item REL: >0.22
            \item $\delta_1$: <0.50
        \end{itemize}
    \end{itemize}

    \item \textbf{Негативный тест: отсутствующие фрагменты}
    \begin{itemize}
        \item \textit{Тестовый набор:} 50 изображений с отсутствующими участками (>30\% площади)
        \item \textit{Ожидаемые метрики:}
        \begin{itemize}
            \item MAE: >0.51
            \item RMSE: >1.09
            \item REL: >0.22
            \item $\delta_1$: <0.50
        \end{itemize}
    \end{itemize}

    \item \textbf{Негативный тест: искаженная геометрия}
    \begin{itemize}
        \item \textit{Тестовый набор:} 50 изображений с сильными деформациями перспективы
        \item \textit{Ожидаемые метрики:}
        \begin{itemize}
            \item MAE: >0.51
            \item RMSE: >1.09
            \item REL: >0.22
            \item $\delta_1$: <0.50
        \end{itemize}
    \end{itemize}
\end{enumerate}

\subsection{Результаты тестирования}

Результаты тестирования представлены в таблице~\ref{tab:test-results-detailed}.

\begin{table}[ht]
\centering
\caption{Результаты тестирования гибридного метода по классам эквивалентности}
\label{tab:test-results-detailed}
\small
\begin{tabular}{|p{5cm}|c|c|c|c|c|}
\hline
\textbf{Класс эквивалентности} & \textbf{MAE} & \textbf{RMSE} & \textbf{REL} & \textbf{$\delta_1$} & \textbf{Вывод} \\
\hline
Оригинальные изображения & 0.312 & 0.621 & 0.138 & 0.832 & Соответствует \\
\hline
Изображения со средними эффектами старения & 0.371 & 0.743 & 0.166 & 0.713 & Соответствует \\
\hline
Изображения с сильными эффектами старения & 0.412 & 0.897 & 0.194 & 0.625 & Соответствует \\
\hline
Негативный тест: экстремально низкое разрешение & 0.638 & 1.245 & 0.276 & 0.412 & Соответствует \\
\hline
Негативный тест: экстремальный уровень шума & 0.715 & 1.518 & 0.318 & 0.357 & Соответствует \\
\hline
Негативный тест: отсутствующие фрагменты & 0.589 & 1.176 & 0.254 & 0.435 & Соответствует \\
\hline
Негативный тест: искаженная геометрия & 0.564 & 1.134 & 0.237 & 0.483 & Соответствует \\
\hline
\end{tabular}
\end{table}

Результаты тестирования подтверждают работоспособность разработанного метода.

\section{Вывод}
В данном разделе был проведен комплексный технологический анализ разработанного гибридного метода восстановления глубины из монокулярных изображений. Определены форматы входных данных (JPEG, PNG, TIFF) и характеристики выходных карт глубины с нормализованными значениями.

Представлены и формализованы ключевые метрики оценки качества реконструкции глубины: MAE, RMSE, REL и точность δ, что позволяет объективно оценивать эффективность метода на разных типах изображений.

Разработана методология тестирования, учитывающая специфику работы с архивными фотографиями через имитацию эффектов старения (сепия, виньетирование, шум, выцветание, дефекты). Использование трех разнородных датасетов (NYU Depth v2, Make3D, SceneNet RGB-D) обеспечило репрезентативность тестирования.

Результаты экспериментов подтверждают устойчивость разработанного метода к различным уровням деградации изображений.
Показатели качества для изображений с сильными эффектами старения (MAE: 0.412, RMSE: 0.897, REL: 0.194, δ₁: 0.625) демонстрируют работоспособность метода даже в сложных условиях, что критически важно для работы с историческими фотографиями.

Проведенные негативные тесты позволили определить граничные условия применимости метода, подтвердив его отказоустойчивость при экстремальных искажениях входных данных.

Ключевой особенностью раздела является разработанная методология тестирования, учитывающая специфику работы с архивными фотографиями через имитацию эффектов старения (сепия, виньетирование, шум, выцветание, дефекты).
Использование трех разнородных датасетов (NYU Depth v2, Make3D, SceneNet RGB-D) обеспечивает репрезентативность тестирования.

Предложенные технологические решения направлены на обеспечение устойчивости метода к различным уровням деградации изображений и определение граничных условий его применимости, что критически важно для работы с историческими фотографиями.



